% sections/limitations_conclusion.tex
% Included from main.tex via \input{sections/limitations_conclusion}

\section{Limitations and Future Work}
\label{sec:limitations}

The current release of CARLA-Air is validated for the workflows presented in
Section~\ref{sec:apps}: single- and dual-drone aerial operations over
moderate-density urban traffic scenes.

\begin{itemize}[leftmargin=*]
  \item \textbf{Actor density.}
  Joint simulation performance is characterized at moderate traffic loads;
  high-density scenes with large simultaneous actor populations remain an
  active engineering target.

  \item \textbf{Environment resets.}
  Map switching requires a full process restart due to independent actor
  lifecycle management in each simulator backend; staged in-session resets are
  planned for a future release.

  \item \textbf{Multi-drone scale.}
  Configurations beyond two drones are functional but not yet formally
  validated across a wide range of scenarios; expanded multi-drone
  characterization will be documented once inter-agent behavior has been fully
  profiled.
\end{itemize}

\noindent None of these constraints affect the workflows in
Section~\ref{sec:apps}, all of which operate within the current boundaries.

Because CARLA-Air inherits and extends AirSim's aerial
capabilities---whose upstream development has been archived---long-term
maintenance of the aerial stack is managed within the CARLA-Air project
itself. Bug fixes, compatibility updates, and feature extensions to the aerial
subsystem are released as part of CARLA-Air's regular update cycle, ensuring
that the flight simulation capabilities continue to evolve independently of
the original upstream repository.

Looking ahead, near-term work will address physics-state synchronization
between the two engines and a ROS\,2~\cite{ros2} bridge that republishes both
simulator streams as standard topics for broader ecosystem integration.
Longer-term, we aim to support GPU-parallel multi-environment execution in the
spirit of Isaac Lab~\cite{isaaclab} and OmniDrones~\cite{omnidrones}, bringing
CARLA-Air closer to the episode throughput required for large-scale
reinforcement learning.


\section{Conclusion}
\label{sec:conclusion}

Simulation platforms for autonomous systems have historically fragmented along
domain boundaries, forcing researchers whose work spans ground and aerial
domains to maintain inter-process bridge infrastructure or accept capability
compromises.
CARLA-Air resolves this fragmentation by integrating CARLA~\cite{carla} and
AirSim~\cite{airsim} within a single Unreal Engine process~\cite{ue4},
exposing both native Python APIs concurrently over a shared physics tick and
rendering pipeline.

\paragraph{Technical contributions.}
The central technical contribution is a principled resolution of the
single-GameMode constraint through a composition-based design:
\texttt{CARLAAirGameMode} inherits CARLA's ground simulation subsystems while
composing AirSim's aerial flight actor as a standard world entity, with
modifications to exactly two upstream source files.
This design yields three properties unavailable in bridge-based approaches: a
shared physics tick that eliminates inter-process clock drift, a shared
rendering pipeline that guarantees consistent weather and lighting across all
sensor viewpoints, and full preservation of both upstream Python APIs.

\paragraph{Platform capabilities.}
Building on this architecture, CARLA-Air provides a photorealistic, physically
coherent simulation world with rule-compliant traffic, socially-aware
pedestrians, and aerodynamically consistent multirotor dynamics.
Up to 18 sensor modalities can be synchronously captured across aerial and
ground platforms at each simulation tick.
The platform directly supports four research directions in air-ground embodied
intelligence: air-ground cooperation, embodied navigation and vision-language
action, multi-modal perception and dataset construction, and
reinforcement-learning-based policy training.

\paragraph{Validation.}
The platform is validated through performance benchmarks demonstrating stable
operation at ${\approx}20$\,FPS under joint workloads, a 3-hour continuous
stability run with zero crashes across 357 reset cycles, and five
representative workflows that exercise the platform's core capabilities.
A consolidated feature comparison with representative platforms is provided in
Table~\ref{tab:platform_comparison} (Section~\ref{sec:related}).

\paragraph{Broader impact.}
By providing a shared world state for aerial and ground agents, CARLA-Air
enables research directions that are structurally inaccessible in single-domain
platforms: paired aerial-ground perception datasets from physically consistent
viewpoints, coordination policies over joint multi-modal observation spaces,
and embodied navigation grounded in cross-view visual and linguistic input.
By also inheriting and extending the aerial simulation capabilities of
AirSim---whose upstream development has been archived---CARLA-Air ensures that
this widely adopted flight stack continues to evolve within a modern, actively
maintained infrastructure.
